\section{自然梯度沿分布几何下降}
\label{sec:12-Real-Analysis-Support/09-Information-Geometry/02}

一个高斯策略网络，输出均值和标准差。把标准差直接当参数学，训练早期动辄发散；改成学 \(\log\sigma\)，同样的学习率、同样的数据、同样的初始化，曲线立刻平顺。模型族一点没变，换的只是记录参数的方式。既然待优化的分布族是同一个，训练轨迹凭什么随记法而变？上一节已经给出答案的前半段：欧氏尺子依赖坐标。这一节把尺子换掉，看下降方向会变成什么。

\subsection{最速下降取决于你允许怎样的一步}

``梯度方向是下降最快的方向''这句话省略了一个前提。把一步更新写成受约束的线性化问题就能看见它：

\[
\Delta^{\star}=\argmin_{\Delta}\ \nabla L(\theta)^{\trans}\Delta
\qquad\text{s.t.}\qquad
\Delta \text{ 落在某个以 } 0 \text{ 为心的小集合内}.
\]

目标函数只用到梯度，所有的选择余地都在那个约束集合的形状里。取欧氏球 \(\norm{\Delta}_2\le\varepsilon\)，解是 \(-\varepsilon\nabla L/\norm{\nabla L}_2\)；取别的形状，解就指向别处。

\begin{center}
\begin{tikzpicture}[>=stealth]
\begin{scope}[shift={(3.0,2.2)}]
% 线性化损失的等值线
\foreach \y in {0.9,0.3,-0.3,-0.9}
  \draw[gray!55,dashed] (-2.3,\y) -- (2.3,\y);
\node[gray!80,right] at (2.4,0.3) {\(\nabla L^{\trans}\Delta=\text{常数}\)};
% 两条支撑线
\draw[gray!70] (-2.3,-1.2) -- (2.3,-1.2);
\draw[gray!70] (-2.3,-1.487) -- (2.3,-1.487);
% 欧氏球与 Fisher 球
\draw[thick] (0,0) circle (1.2);
\draw[thick,rotate around={40:(0,0)}] (0,0) ellipse (2.2 and 0.6);
% 两个方向
\draw[->,very thick] (0,0) -- (0,-1.2);
\draw[->,very thick] (0,0) -- (-1.483,-1.487);
\fill (0,0) circle (0.06);
% 标注
\draw[gray!70] (-1.05,1.75) -- (-0.6,1.039);
\node[left] at (-1.05,1.75) {\(\norm{\Delta}_2\le\varepsilon\)};
\draw[gray!70] (1.05,1.75) -- (1.40,1.50);
\node[right] at (1.05,1.75) {\(\Delta^{\trans}I(\theta)\Delta\le\varepsilon^{2}\)};
\node[fill=white,inner sep=1pt,right] at (0.12,-1.32) {\(-\nabla L\)};
\node[below] at (-1.483,-1.62) {\(-I(\theta)^{-1}\nabla L\)};
\end{scope}
\end{tikzpicture}
\end{center}

\emph{同一个梯度、同一组线性化等值线；信任域是圆时最速方向朝正下方，信任域是 Fisher 椭圆时它偏向椭圆被拉长的那一侧。}

图里两个箭头都指向各自信任域上``线性化损失最低''的那个边界点。圆是各向同性的，所以答案与 \(\nabla L\) 平行；椭圆在某个方向上宽，允许在那个方向上多走，答案随之偏转。把约束换成 Fisher 椭圆

\[
\Delta^{\trans}I(\theta)\,\Delta\le\varepsilon^2,
\]

意思是限制这一步造成的\emph{分布}改变，而不限制参数数值的改变——因为上一节已经说明，这个二次型度量的正是分布的可区分程度。

\subsection{自然梯度是 Fisher 度量下的那个代表元}

同样的东西可以用内积的语言说一遍，得到的公式更好用。方向导数 \(v\mapsto \nabla L(\theta)^{\trans}v\) 是切空间上的一个线性泛函。所谓梯度，是把这个泛函用内积表示出来的那个向量：在欧氏内积下它是 \(\nabla L\)，在 Fisher 内积 \(g_\theta(u,v)=u^{\trans}I(\theta)v\) 下，要找的向量 \(\widetilde\nabla L\) 满足

\[
v^{\trans}\nabla L(\theta)=g_\theta\bigl(\widetilde\nabla L(\theta),\,v\bigr)
=v^{\trans}I(\theta)\,\widetilde\nabla L(\theta)
\qquad\text{对一切 } v .
\]

两边对任意 \(v\) 相等，于是

\[
\widetilde\nabla L(\theta)=I(\theta)^{-1}\nabla L(\theta).
\]

这个向量叫自然梯度，对应的更新是 \(\theta_{k+1}=\theta_k-\alpha I(\theta_k)^{-1}\nabla L(\theta_k)\)。它与信任域的说法给出同一个方向，只差一个正的缩放因子。

\subsection{换坐标时自然梯度跟着走，欧氏梯度不跟}

作可逆光滑变换 \(\eta=\eta(\theta)\)，记 \(J=\partial\eta/\partial\theta\)。链式法则给出 \(\nabla_\theta L=J^{\trans}\nabla_\eta L\)，而上一节的度量变换律是 \(I_\theta=J^{\trans}I_\eta J\)。两者一合：

\[
\widetilde\nabla_\theta L
=I_\theta^{-1}\nabla_\theta L
=J^{-1}I_\eta^{-1}J^{-\trans}J^{\trans}\nabla_\eta L
=J^{-1}\,\widetilde\nabla_\eta L .
\]

这正是切向量在坐标变换下的规律。翻译成实际操作：在 \(\theta\) 坐标下走一步自然梯度，与先换到 \(\eta\) 坐标、在那里走一步自然梯度、再换回来，得到同一个点（就一阶而言）。欧氏梯度做不到这一点，它按 \(J^{\trans}\) 变换，方向会随坐标漂移。开头那个 \(\sigma\) 与 \(\log\sigma\) 的差别，在自然梯度下消失。

这种``对坐标变换免疫''的性质并不新鲜。Newton 法的仿射不变性说的是同一件事：用 Hessian 的逆作预条件，问题作可逆线性变换后迭代轨迹不变，见\hyperref[sec:09-Convex-Optimization/06-Optimization-Algorithms/02]{``Newton 法用曲率校正方向''}。两者的形式都是``正定矩阵的逆乘梯度''，区别在矩阵的来源：Hessian 来自目标函数在当前点的曲率，会随数据批次波动、也可能不正定；Fisher 信息只依赖模型族 \(\set{p_\theta}\)，按定义半正定。极大似然的情形下二者渐近一致，因为对数似然的 Hessian 期望等于 Fisher 信息的相反数。

\subsection{Bernoulli 上一算，两处发散互相抵消}

\(n\) 次独立试验成功 \(S\) 次，负对数似然 \(L(\theta)=-S\log\theta-(n-S)\log(1-\theta)\)，欧氏梯度

\[
\nabla_\theta L=-\frac{S}{\theta}+\frac{n-S}{1-\theta}
=\frac{n\theta-S}{\theta(1-\theta)}
\]

在 \(\theta\to0\) 处发散。Fisher 信息是 \(n/\bigl(\theta(1-\theta)\bigr)\)，它同样发散，而它的逆趋于零。两处发散相乘：

\[
\widetilde\nabla_\theta L
=\frac{\theta(1-\theta)}{n}\cdot\frac{n\theta-S}{\theta(1-\theta)}
=\frac{n\theta-S}{n}
=\theta-\widehat\theta ,
\]

其中 \(\widehat\theta=S/n\)。下降方向 \(-\widetilde\nabla L=\widehat\theta-\theta\) 始终直指最大似然解，步长 \(\alpha=1\) 一步到位。边界附近欧氏梯度的爆炸，被度量的逆恰好中和；这两处的发散源出同一个 \(\theta(1-\theta)\)，抵消不是巧合。

\subsection{实践中拿到手的都是近似，代价各不相同}

理想的自然梯度在深度学习里没法直接算。参数量 \(d\) 达到 \(10^8\) 以上时，\(I(\theta)\) 有 \(d^2\) 个元素，求逆是 \(O(d^3)\)，光是存下来就不可能。落地的做法都是在近似 \(I(\theta)\)，代价体现在两处：近似得多粗，以及坐标不变性还剩多少。

最省事的是只保留对角元。这时预条件矩阵是 \(\diag\bigl(I_{11},\dots,I_{dd}\bigr)\)，求逆退化成逐坐标除法。Adam 走的接近这条路线：它按坐标累积梯度平方的滑动平均 \(\hat v_t\)，再用 \(\hat v_t^{-1/2}\) 缩放梯度。梯度平方的期望在对数似然损失下正是经验 Fisher 的对角元，所以 Adam 可以读成一个很粗的对角自然梯度——只是它开了平方根，得到的是 \(I^{-1/2}\) 而非 \(I^{-1}\) 量级的缩放。这一点让它在实践中更保守，也让它与自然梯度的对应关系只是形似。对角近似换来的不变性也相应缩水：它对每个坐标各自的缩放免疫，对混合了不同坐标的一般变换无能为力。参数间强相关时，被忽略的非对角块正是病态的来源，可参考\hyperref[sec:14-Deep-Learning/03-Optimization-and-Training/02]{``动量与 Adam 在平均什么''}。

K-FAC 走得更远一点。它按层拆开，把每层的 Fisher 块近似成两个小矩阵的 Kronecker 积 \(A\otimes G\)，其中 \(A\) 来自该层输入的二阶统计、\(G\) 来自反传信号的二阶统计。Kronecker 积的逆等于两个因子的逆的 Kronecker 积，于是 \(O(d^3)\) 被压到两个小矩阵求逆。代价是那条``输入统计与梯度统计独立''的假设并不成立，以及要周期性重估并缓存 \(A,G\) 的逆，显存与调度都要额外安排。

还有两处细节在工程上绕不开。其一是用经验 Fisher（在观测标签上平均）代替真 Fisher（在模型自己的预测分布上取期望），二者在模型拟合不好时差别明显。其二是阻尼：\(I(\theta)\) 在过参数化模型上接近奇异，实际用的总是 \(I(\theta)+\lambda \mathrm{Id}\)，而 \(\lambda\) 一旦调大，更新就退回普通梯度下降。回看那张图，阻尼做的事情是把 Fisher 椭圆往圆的方向捏。

\subsection{它改的是方向，不是问题本身}

自然梯度不会让非凸目标变凸，不会消除随机梯度的噪声，也修不了设定错误的模型。它只回答一个问题：给定要下降的目标和要保护的分布，这一步该往哪儿迈。

约束``一步造成的分布改变不超过多少''这个想法本身，比自然梯度的具体公式走得更远——强化学习里的信任域方法直接把 KL 散度写进约束，而不是用它的二阶近似。要理解那条约束的几何含义，还差最后一块拼图：KL 散度与 Fisher 度量在局部到底是什么关系。下一节补上。
